% concatenated from App.tex
 %\documentclass[11pt]{article}

%\usepackage{amsmath, amssymb, amsthm}
%\usepackage{mathrsfs}
%\usepackage[margin=1in]{geometry}

% Standard notation
\newcommand{\QQ}{\mathbf{Q}}
\newcommand{\ZZ}{\mathbf{Z}}
\newcommand{\FF}{\mathbf{F}}
\newcommand{\HH}{\mathbf{H}}
\newcommand{\RR}{\mathbf{R}}
\newcommand{\CC}{\mathbf{C}}
\newcommand{\PP}{\mathbf{P}}

\renewcommand{\AA}{\mathbf{A}}

\newcommand{\abs}[1]{\lvert\!\lvert #1 \rvert\!\rvert}
\newcommand{\zmod}[1]{\,\,(\mathrm{mod}\,\,#1)}
\newcommand{\Gal}{\mathrm{Gal}}
\newcommand{\pp}{\mathfrak{p}}
\newcommand{\qq}{\mathfrak{q}}
\newcommand{\ga}{J_{\pi}}
\newcommand{\gb}{J_{\pi'}}
\newcommand{\onto}{\twoheadrightarrow}

% References used in the text
\newcommand{\Sref}{1}
\newcommand{\Zref}{2}

% Theorem environments
\newtheorem{proposition}{Proposition}
\renewcommand{\theproposition}{B.\arabic{proposition}}

\newtheorem{corollary}[proposition]{Corollary}

%\title{Appendix to ``Maximal curves of genus $5$ over finite fields''}
%\author{
%Ren\'e Schoof\\[0.5em]
%\small Dipartimento di Matematica\\
%\small $2^{\mathrm{a}}$ Universit\`a di Roma ``Tor Vergata''\\
%\small I-00133 Roma, Italy\\
%\small \texttt{schoof@mat.uniroma2.it}
%}
%\date{}

%\begin{document}

%\maketitle

\section*{Appendix B to ``Maximal curves of genus $5$ over finite fields"} \label{appendix B}

Let $X$ be a smooth, projective, absolutely irreducible curve over a finite field
$\FF_q$ of genus $g_X=5$. Suppose that $X$ admits an automorphism of order $5$
and let $G$ denote the group generated by it. Let $E$ denote the quotient of $X$
by $G$. The conductor-discriminant formula applied to the cyclic Galois cover
$X \rightarrow E$ is
\[2g_X - 2 = 5(2g_E - 2)+\sum_{\chi} \deg(\mathrm{conductor}(\chi)).
\]
It easily implies that the genus $g_E$ of $E$ is $1$ and that the degree of the
common conductor $D$ of the non-trivial characters $\chi$ of $G$ is equal to $2$.

By class field theory the reciprocity homomorphism
\[
\theta : \AA_E^*/K^*U_D \longrightarrow G
\]
is surjective. Here $K$ is the function field $\FF_q(E)$ and $\AA_E$ is the
ad\`ele ring of $E$. The subgroup of the id\`ele group $\AA_E^*$ of unit
id\`eles is denoted by $U$. The quotient group $\AA_E^*/U$ is naturally
isomorphic to the divisor group $\mathrm{Div}(E)$. Let $U_D$ denote the subgroup
of $u \in U$ for which $u \equiv 1 \zmod D$. There is an exact sequence
\begin{equation*}
 \begin{matrix}
0 & \longrightarrow & U/\FF_q^*U_D
& \longrightarrow & \AA_E^*/K^*U_D
& \longrightarrow & \AA_E^*/K^*U
& \longrightarrow & 0 \\[0.3em]
&&&&&& \Vert && \\[0.3em]
&&&&&& \mathrm{Pic}(E) &&
\end{matrix}   \tag{$\star$}
\end{equation*}

 
Since the degree of $D$ is $2$, there are three possibilities for $D$. It is
either a point of degree $2$, or it is the sum of two degree $1$ points, which
may or may not be equal. The group $U/\FF_q^*U_D$ is isomorphic to
$\FF_{q^2}^*/\FF_q^*$, $\FF_q^*$ or $\FF_q$ in these cases. The orders of these
groups are $q+1$, $q-1$ and $q$, respectively.

\begin{proposition}
We have $q \not\equiv \pm 2 \zmod 5$.
\end{proposition}
\begin{proof}
If $q$ were congruent to $\pm 2 \zmod 5$, none of the possible orders of
$U/\FF_q^*U_D$ is divisible by $5$. Therefore $U/\FF_q^*U_D$ is in the kernel of
$\theta$. It follows that $\theta$ factors through
$\AA_E^*/K^*U = \mathrm{Pic}(E)$. But that implies that the cover
$X \rightarrow E$ is unramified, which it is not. This contradiction proves the
proposition.
\end{proof}

By Zaytsev~\cite[Thm.~4.13]{zaytsevoptimal}, an optimal genus $5$ curve over $\FF_q$
for which one has $[2\sqrt{q}]^2 - 4q = -19$, admits an automorphism of order $5$. Therefore Proposition~B.1 applies. It gives an alternative proof of Theorem \ref{Theorem 3.2} of this paper.

\vspace{0.5cm}

The next proposition regards $q=5$. Over $\FF_5$ there is a unique elliptic
curve $E'$ with $\#E'(\FF_5)=5$. The trace of its Frobenius endomorphism is $1$
and the discriminant of its endomorphism ring is $-19$. The curve $E'$ is given
by the Weierstrass equation $y^2 = x^3 - 2x + 2$.


\begin{proposition}
If $q=5$, then the curve $E$ cannot be isomorphic to $E'$.
\end{proposition}

\begin{proof}
If $q=5$, the degree $2$ divisor $D$ can only be $2P$ for some
$\FF_q$-rational point $P$ of $E$. In this case the order of
$U/\FF_q^*U_D$ is equal to $q=5$. The exact sequence $(\star)$ modulo fifth powers
leads to the exact sequence
\[
U/\FF_5^*U_D
\longrightarrow
\AA_E^*/{\AA_E^*}^5K^*U_D
\longrightarrow
\AA_E^*/{\AA_E^*}^5K^*U
\longrightarrow 0.
\]
If $E$ were isomorphic to the elliptic curve $E'$ given by
$y^2=x^3-2x+2$, the leftmost homomorphism is zero. In other words, $U$ is
contained in the subgroup ${\AA_E^*}^5K^*U_D$ of $\AA_E^*$. To see this, we first observe that the group $E(\FF_5)$ acts transitively on
itself by translations. Therefore we may take for $P$ any of the five points in
$E(\FF_5)$. We choose $P=(1,1)$. Then we put $Q=(2,-1)$ and let
$\xi \in \AA_E^*$ denote an id\`ele whose image in $\mathrm{Div}(E)$ is the
divisor $Q-O$, where $O$ is the point of $E$ at infinity. It is not difficult to see that the divisor of the function
\[
h =
\frac{(y+1)^2(x-2)}{y-2x-2}
\]
is equal to $5(Q-O)$. Therefore both $\xi^5$ and $h$ are elements of
$\AA_E^*$ whose images in $\mathrm{Div}(E)$ are equal to $5(Q-O)$. It follows
that $v = \frac{\xi^5}{h}$ is in $U$. Since $h(P)=2$ and
\[
h-2 =
\frac{(x-1)(2y+x^3-x^2+2x)}{y-2x-2}
\]
has a simple zero in $P$, the id\`ele unit $v$ generates $U/U_D$. Therefore
every $u \in U$ can be written as $v^k w$ for some $k \in \ZZ$ and
$w \in U_D$. Since $v$ is in ${\AA_E^*}^5K^*$, this proves that we indeed have $U \subset {\AA_E^*}^5K^*U_D$.
 

It follows that the natural map
\[
\AA_E^*/{\AA_E^*}^5K^*U_D
\mathop{\longrightarrow}^{\cong}
\AA_E^*/{\AA_E^*}^5K^*U
\]
is an isomorphism. By class field theory, the group on the left is isomorphic
to the Galois group of the maximal abelian exponent $5$ extension $L$ of $K$ of
conductor at most $D=2P$. Note that the function field $\FF_5(X)$ of the curve
$X$ is a subfield of $L$. Since $\AA_E^*/{\AA_E^*}^5K^*U_D$ is isomorphic to  $\AA_E^*/{\AA_E^*}^5K^*U$,
 the extension $K \subset L$ is unramified at \textit{all} places. Since we have $K \subset \FF_5(X) \subset L$,
 the same is true for $\FF_5(X)$. However, the cover $X \rightarrow E'$ is
actually ramified and hence we obtain a contradiction.
This proves the proposition.
\end{proof}

\begin{corollary}
When $q$ is a power of $5$, there is no optimal genus $5$ curve $X$ over
$\FF_q$ for which $[2\sqrt{q}]^2 - 4q = -19.$
\end{corollary}

\begin{proof}
Suppose $X$ is such an optimal curve. By Proposition \ref{Diophantine} of this paper, the
curve $X$ cannot exist, except possibly when $q=5^7$. In this case we have $[2\sqrt{q}] = 559$ and the eigenvalues of the characteristic polynomial of Frobenius are $\frac{-559+\sqrt{-19}}{2} = \phi^7$ and its complex conjugate, where $\phi = \frac{1+\sqrt{-19}}{2}$. Since the ring $\ZZ[\phi]$ is generated by $\phi^7$, a theorem of Serre~\cite[Thm.~9 of the appendix]{Lauter} implies that $X$ is the base change
of a genus $5$ curve $X'$ over $\FF_5$.

By Zaytsev~\cite[Thm.~4.13]{zaytsevoptimal}, the automorphism group of $X$ over
$\FF_{5^7}$ is isomorphic to the dihedral group $D_5$. The Galois group of
$\FF_{5^7}$ over $\FF_5$ acts on $\mathrm{Aut}_{\FF_{5^7}}(X')$. Since $7$ is
prime to $\#\mathrm{Aut}(D_5)=20$, this action is trivial. Therefore all
automorphisms of $X'$ are defined over $\FF_5$. The quotient of $X'$ by an
automorphism of order $5$ is isomorphic to the unique elliptic curve $E'$ over
$\FF_5$ with five rational points. Therefore Proposition B.2 applies and we
conclude that this is not possible. Therefore $X$ cannot exist.
\end{proof}

In view of the results in this paper, the fields $\FF_q$ over which an optimal
genus $5$ curve $X$ with $[2\sqrt{q}]^2 - 4q = -19$ may exist necessarily satisfy $q \equiv 1 \zmod 5$. As we explained above, in these cases $X$ is a cyclic degree $5$ cover ramified
over two distinct $\FF_q$-rational points of a specific genus $1$ curve $E$. By
translating we may assume that the two points are the point $O$ at infinity and
a second point $P$. By Kummer theory the function field of $X$ is of the form $\FF_q(E)\left(\sqrt[5]{h}\right)$, where $h \in \FF_q(E)$ is a function whose divisor is $m(P-O)$ with
$m=\#E(\FF_q)$. The function $h$ is unique up to a scalar multiple. The number
of $\FF_q$-rational points on $X$ is equal to $2+5r$, where $r$ is the number
of points $Q \in E(\FF_q) - \{P,O\}$ for which $h(Q)$ is a fifth power in $\FF_q^*$.

A short Pari-GP program, computing the number $r$ for each point $P \neq O$ in
$E(\FF_q)$ and each function $h$, showed that for $q<10000$ there is no optimal
genus $5$ curve $X$ over $\FF_q$ with $[2\sqrt{q}]^2 - 4q = -19$. This is actually a short list of ten primes, the smallest and largest being
$q=61$ and $q=9511$, respectively. The entire computation took less than two
hours.

 

 

\vspace{1em}

\noindent
Ren\'e Schoof\\[0.5em]
Dipartimento di Matematica\\
$2^{\mathrm{a}}$ Universit\`a di Roma ``Tor Vergata''\\
I-00133 Roma, Italy\\
Email: \texttt{schoof@mat.uniroma2.it}

%\end{document}
